% Draft prose bounded by verified artifacts; replace citation keys with bibliography entries.
\section{Methods}
\paragraph{Datasets and study separation.} We analyzed N-MNIST in two distinct studies. In the QSNN-v3 timestamp study, models were trained and selected using a frozen 55,000/5,000 split of the official 60,000-sample training partition; the official test partition was not evaluated for QSNN-v3. The attack set comprised 100 seed-42 validation samples (10 per class) correctly classified by both the compact SNN and QSNN-v3. In the controlled architecture study, four models were retrained on a common stratified 55,000/5,000 split, evaluated on all 10,000 official test samples, and attacked on the first 1,000 test samples jointly classified correctly by all four. Binary and integer count representations used separate equal-event-count, ten-bin caches and frozen manifests. These studies have distinct representations and threat models and are reported separately. [Sources: `scripts/run_nmnist_common_attack_protocol_seed42.py:100-123`; `configs/nmnist_controlled_four_models_lr1e4.json`; `scripts/build_nmnist_controlled_manifest.py`.]

\paragraph{Models and training.} The compact SNN comprises two Conv/BatchNorm/average-pooling LIF stages and a ten-class linear head over time-averaged second-stage spikes, with 25,482 trainable parameters. QSNN-v3 uses a 32-dimensional Conv/LIF feature extractor, learned eight-angle projection, eight-qubit two-block variational state-vector circuit, learned measurement basis and 256-probability linear classifier without a classical bypass (40,762 parameters). Its five-seed clean study used AdamW with group learning rates and validation-selected checkpoints; the controlled four-model study used a common Adam learning rate $10^{-4}$, batch size 64 and the same validation-selection rule for the compact SNN, ConvNet, Spiking ResNet18 and VGGSNN. [Sources: `models/nmnist_snn.py:23-55`; `models/nmnist_hybrid_qsnn.py:59-279`; `scripts/run_nmnist_qsnn_v3_multiseed.py:52-100`; `configs/nmnist_controlled_four_models_lr1e4.json`.]

\paragraph{Timestamp attacks.} For seed-42 QSNN-v3 and compact SNN checkpoints, untargeted white-box PGD maximized true-label cross entropy with 20 signed timestamp steps, while TEMP-DRIFT-v2 used derivative-free timestamp search. Both preserved event count, coordinates, polarity and temporal order, with each timestamp bounded by $\epsilon=\rho(t_N-t_1+1)$ around the clean timestamp. We compared access-matched and independently calibrated wall-clock-matched conditions on the same clean-correct samples, recording attack forwards, backward evaluations, per-sample success and distortion. [Sources: `scripts/run_nmnist_attack_protocol_v3_auditable_seed42.py:238-253,304-345`; `scripts/run_nmnist_budget_matched_comparison_v2_seed42.py:686-695`.]

\paragraph{Grid-cell attacks.} In the separate controlled four-model study, PIL-PGD acted on nonzero temporal-grid cells as indivisible packets; integer amplitudes remained attached to their packets. Strict projection preserved packet count, amplitude, polarity, spatial line and collision-free occupancy. We evaluated budgets $B_\infty=1$ and $B_0=200$ at $T=10$ on the same 1,000 jointly clean-correct samples per representation and computed untargeted ASR with that fixed denominator. Serialized adversarial examples were subjected to independent reconstruction, budget, prediction and hash audits. [Sources: `AGENTS.md`; `scripts/run_nmnist_controlled_attack.py:25-133`; `scripts/audit_nmnist_controlled_attack.py`.]

\paragraph{Analysis boundary.} All robustness comparisons are seed-42 observations. Published reference-paper values and earlier locally trained custom-SNN rows are descriptive context, not interchangeable with the common-recipe four-model measurements. We do not infer defended-model performance, QSNN-v3 official-test accuracy or multi-seed robustness from these experiments. [Sources: `readme_jafar.md:119-159,175-190,282-286`; `Reports/nmnist_controlled_four_models_lr1e4_comparison.md`.]
